\documentclass[14pt,dvipdfmx]{jsarticle}
\input{lecture-preamble.tex}

\setcounter{definition}{2}
\setcounter{theorem}{7}
\setcounter{technique}{0}

\begin{document}
\lectureheader{数学Ⅱ}{5}{指数関数と対数関数}{2}{対数関数}{4-C}{対数関数を含む方程式, 不等式}{167--168}{対数関数を含む方程式・不等式について考えてみよう}

% ============ 1ページ目（表・左） ============
\vspace{-1.8zw}
\heading{〜前時の復習〜}
\vspace{-0.55em}
\begin{reviewbox}{対数関数とそのグラフ}
  \begin{enumerate}[label=\textbf{\arabic*}, leftmargin=1.5zw, itemsep=0.02em]
    \item $y=\log_{a} x$ のグラフは $y=a^{x}$ と直線 $y=\fitblankbf{x}$ に関して対称.
    \item 点 $(1,\ \fitblankbf{0})$, $(a,\ \fitblankbf{1})$ を通り, \fitblankbf{y \text{軸}}を漸近線とする.
    \item $a>1$ のとき\fitblankbf{\text{増加関数}}, $0<a<1$ のとき\fitblankbf{\text{減少関数}}.
  \end{enumerate}
\end{reviewbox}

\heading{対数関数を含む方程式}
\vspace{-0.25em}
対数方程式も指数方程式と同様に, 今まで学習してきた知識を使って \fitblankbf{x = \text{○}} の形に変形する.

\begin{exampleprob}{6}
  次の方程式を解け。\\
  (1)\quad $\log_{2} x = 3$
\end{exampleprob}
\begin{answerbox}{(1)}[3.8]
  対数の定義から \qquad $x=2^{3}=8$
\end{answerbox}

\vspace{2em}
\begin{techniquebox}{対数方程式の解き方}
  \begin{enumerate}[label=\textbf{\arabic*}, leftmargin=1.5zw, itemsep=0.02em]
    \item 真数が\fitblankbf{\text{正}}である範囲を求める.
    \item 対数の定義 $(\log_{a} x = p \Longleftrightarrow x = \fitblankbf{a^{p}})$ や

    対数の性質$(\log_{a} M+\log_{a} N=\fitblankbf{\log_{a} MN}など)$を使って式を変形する。
    \item 求めた解が, \fitblankbf{真数条件}をみたすか確かめる.
  \end{enumerate}
\end{techniquebox}

\vfill
\nextpage

% ============ 2ページ目（表・右） ============
\begin{practiceprob}{21}
  次の方程式を解け。
  \begin{tasks}[label=(\arabic*), label-width=1.6zw, item-indent=2zw, after-item-skip=0.08em](2)
    \task $\log_{3} x = 4$
    \task $\log_{\frac{1}{2}} x = -2$
  \end{tasks}
\end{practiceprob}
\enlargethispage{6zw}
\begin{minipage}[t]{0.48\textwidth}
\begin{answerbox}{(1)}[4.5]
  対数の定義から

  $x=3^{4}=81$
\end{answerbox}
\end{minipage}\hfill
\begin{minipage}[t]{0.48\textwidth}
\begin{answerbox}{(2)}[4.5]
  対数の定義から

  $x=2^{2}=4$
\end{answerbox}
\end{minipage}

\noindent{\bfseries 応用例題\,1}\quad
方程式 $\log_{3} x+\log_{3}(x-8)=2$ を解け。
\vspace{0.15em}

\begin{answerbox}{}[15]
  \makebox[10zw][l]{真数が正であるから}$x>0$ かつ $x>8$ \\
  \makebox[10zw][l]{すなわち}$x>8$ ……①\\
  \makebox[10zw][l]{方程式を変形すると}$\log_{3} x(x-8)=2$\\
  \makebox[10zw][l]{よって}$x(x-8)=3^{2}$\\
  $(x+1)(x-9)=0$ より, ①をみたすのは $x=9$
\end{answerbox}

\begin{practiceprob}{22}
  次の方程式を解け。
  \begin{tasks}[label=(\arabic*), label-width=1.6zw, item-indent=2zw, after-item-skip=0.06em](2)
    \task $\log_{2} x+\log_{2}(x-6)=3$
    \task $\log_{2}(x+1)+\log_{2}(x-2)=2$
  \end{tasks}
\end{practiceprob}
\begin{minipage}[t]{0.48\textwidth}
\begin{answerbox}{(1)}[8.4]
  $x>6$ ……①\\
  $x(x-6)=8$,\ $(x-8)(x+2)=0$\\
  ①より $x=8$
\end{answerbox}
\end{minipage}\hfill
\begin{minipage}[t]{0.48\textwidth}
\begin{answerbox}{(2)}[8.4]
  $x>2$ ……①\\
  $(x+1)(x-2)=4$,\ $(x-3)(x+2)=0$\\
  ①より $x=3$
\end{answerbox}
\end{minipage}

\vspace{0.2em}
% \noindent{\bfseries 演習問題}\quad
% 次の方程式を解け。
% \begin{tasks}[label=(\arabic*), label-width=1.6zw, item-indent=2zw, after-item-skip=0.06em](2)
%   \task $\log_{6} x+\log_{6}(x+1)=1$
%   \task $2\log_{2}(x-3)=\log_{2}(5-x)$
% \end{tasks}
% \vspace{0.08em}
% \begin{minipage}[t]{0.48\textwidth}
% \begin{answerbox}{(1)}[8.6]
%   $x>0$ ……①\\
%   $x(x+1)=6$,\ $(x+3)(x-2)=0$\\
%   ①より $x=2$
% \end{answerbox}
% \end{minipage}\hfill
% \begin{minipage}[t]{0.48\textwidth}
% \begin{answerbox}{(2)}[8.6]
%   $3<x<5$ ……①\\
%   $(x-3)^{2}=5-x$,\ $(x-1)(x-4)=0$\\
%   ①より $x=4$
% \end{answerbox}
% \end{minipage}

\vfill
\nextpage

% ============ 3ページ目（裏・左） ============
\heading{対数関数を含む不等式}
\vspace{-0.1em}

\begin{exampleprob}{6}
  (2)\quad $\log_{2} x \leqq 3$
\end{exampleprob}
\begin{answerbox}{(2)}[6.6]
  真数が正であるから $x>0$ ……①\\
  $\log_{2} x \leqq \log_{2} 8$, 底 $2>1$ より $x\leqq 8$ ……②\\
  ①, ②より $0<x\leqq 8$
\end{answerbox}

\begin{techniquebox}{対数不等式の解き方}
  \begin{enumerate}[label=\textbf{\arabic*}, leftmargin=1.5zw, itemsep=0.02em]
    \item 真数が\fitblankbf{\text{正}}である範囲を先に求める.
    \item 両辺を同じ底の対数にする.
    \item 底 $a>1$ なら向きは\fitblankbf{\text{そのまま}}, $0<a<1$ なら向きは\fitblankbf{\text{逆}}.
    \item 真数の条件との\fitblankbf{\text{共通範囲}}を求める.
  \end{enumerate}
\end{techniquebox}

\begin{practiceprob}{21}
  次の不等式を解け。
  \begin{tasks}[label=(\arabic*), label-width=1.6zw, item-indent=2zw, after-item-skip=0.06em](2)
    \task[(3)] $\log_{2} x \leqq 4$
    \task[(4)] $\log_{\frac{1}{2}} x < 2$
    \task[(5)] $\log_{3} x \geqq \dfrac{1}{2}$
    \task[(6)] $\log_{\frac{1}{3}} x > -1$
  \end{tasks}
\end{practiceprob}
\vspace{0.08em}
\begin{minipage}[t]{0.48\textwidth}
\begin{answerbox}{(3)}[6.6]
  $x>0$ ……①\\
  底 $2>1$ より $x\leqq 16$ ……②\\
  ①, ②より $0<x\leqq 16$
\end{answerbox}
\end{minipage}\hfill
\begin{minipage}[t]{0.48\textwidth}
\begin{answerbox}{(4)}[6.6]
  $x>0$ ……①\\
  底 $\dfrac{1}{2}<1$ より $x>\dfrac{1}{4}$ ……②\\
  ①, ②より $x>\dfrac{1}{4}$
\end{answerbox}
\end{minipage}
\vspace{0.12em}
\begin{minipage}[t]{0.48\textwidth}
\begin{answerbox}{(5)}[6.6]
  $x>0$ ……①\\
  底 $3>1$ より $x\geqq \sqrt{3}$ ……②\\
  ①, ②より $x\geqq \sqrt{3}$
\end{answerbox}
\end{minipage}\hfill
\begin{minipage}[t]{0.48\textwidth}
\begin{answerbox}{(6)}[6.6]
  $x>0$ ……①\\
  底 $\dfrac{1}{3}<1$ より $x<3$ ……②\\
  ①, ②より $0<x<3$
\end{answerbox}
\end{minipage}

\vfill
\nextpage

% ============ 4ページ目（裏・右） ============
\noindent{\bfseries 応用例題\,2}\quad
不等式 $\log_{\frac{1}{3}}(x-1)>2$ を解け。
\vspace{0.08em}

\begin{answerbox}{解}[6.6]
  真数が正であるから $x>1$ ……①\\
  底 $\frac{1}{3}<1$ より $x-1<\frac{1}{9}$ ……②\\
  ①, ②より $1<x<\frac{10}{9}$
\end{answerbox}

\begin{practiceprob}{23}
  次の不等式を解け。
  \begin{tasks}[label=(\arabic*), label-width=1.6zw, item-indent=2zw, after-item-skip=0.04em](2)
    \task $\log_{2}(x-3)<4$
    \task $\log_{\frac{1}{2}}(x+2)\geqq 3$
  \end{tasks}
\end{practiceprob}
\begin{minipage}[t]{0.48\textwidth}
\begin{answerbox}{(1)}[6.4]
  $x>3$ ……①,\ $x-3<16$ ……②\\
  ①, ②より $3<x<19$
\end{answerbox}
\end{minipage}\hfill
\begin{minipage}[t]{0.48\textwidth}
\begin{answerbox}{(2)}[6.4]
  $x>-2$ ……①,\ $x+2\leqq\frac{1}{8}$ ……②\\
  ①, ②より $-2<x\leqq -\frac{15}{8}$
\end{answerbox}
\end{minipage}

\vspace{0.15em}
% \noindent{\bfseries 演習問題}\quad
% 次の不等式を解け。
% \begin{tasks}[label=(\arabic*), label-width=1.6zw, item-indent=2zw, after-item-skip=0.04em](2)
%   \task $\log_{2}(x+1)>\log_{2}(2x-5)$
%   \task $2\log_{\frac{1}{3}}(x-1)<\log_{\frac{1}{3}}(13-x)$
% \end{tasks}
% \vspace{0.06em}
% \begin{minipage}[t]{0.48\textwidth}
% \begin{answerbox}{(1)}[7.2]
%   $x>\frac{5}{2}$ ……①\\
%   底 $2>1$ より $x<6$ ……②\\
%   ①, ②より $\frac{5}{2}<x<6$
% \end{answerbox}
% \end{minipage}\hfill
% \begin{minipage}[t]{0.48\textwidth}
% \begin{answerbox}{(2)}[7.2]
%   $1<x<13$ ……①\\
%   $(x-4)(x+3)>0$\\
%   ①より $4<x<13$
% \end{answerbox}
% \end{minipage}

\vspace{0.15em}
\begin{summarybox}{対数関数を含む方程式・不等式}
  \begin{enumerate}[label=\textbf{\arabic*}, leftmargin=1.5zw, itemsep=0.06em]
    \item 方程式 $\log_{a} x = p$ は, 定義から $x=\fitblankbf{a^{p}}$ と解く. 真数は\fitblankbf{\text{正}}.
    \item 不等式では, 真数条件を先にとり, 底 $a>1$ なら向きは\fitblankbf{\text{そのまま}},\\
          $0<a<1$ なら向きを\fitblankbf{\text{逆}}にする.
  \end{enumerate}
\end{summarybox}

\vfill
\nexttime{常用対数を使って, 桁数を求めてみよう!}
\end{document}
